
\subsection{Zigzags in a range category \catcw}

For any $n \ge 1$, a zigzag of length $n$ in range category \catcw  
is a sequence of morphisms  $F=\zigzagTuple{f}{n}$ of \catcw such that for some 
objects $x_0,..x_n$ and $y_1,...y_n$ of \catcw we have
$$
\arraycolsep=7pt
\zigzag{x}{y}{f}{n}
$$
in \catcw. 

Within a zigzag of length $n$ we find a zig for each \foreachi[n-1], like so
$$
\begin{array} {c c c}
&\Ob{y}{i}&\\[0.8cm]
\Ob{x}{i}&&\Obs{x}{i}
\end{array}
\begin{arrows}
\ncarr{xi}{yi}   \alabel{f_i}[0.45][0.4]
\ncarr {xis}{yi} \blabel{f'_i}[0.45][0.4]
\end{arrows}
$$

We will say that this i'th zig is \textit{tight} iff
$\rng{f_i}=\rng{f'_i}$.

We also find a  zag for each $i$, $2 \leq i \leq n$, like so
%\arreset}
$$
\setlength{\arrnodesepA}{3pt}
\begin{array} {c c c}
\Obp{y}{i}&&\Ob{y}{i} \\[0.8cm]
&\Ob{x}{i}&
\end{array}
\begin{arrows}
\ncarr{xi}{yip}  \alabel{f'_i}[0.45][0.4]
\ncarr {xi}{yi}  \blabel{f_i}[0.45][0.4]
\end{arrows}
$$

and we will say that the i'th zag is tight iff
$\rst{f'_i}=\rst{f_i}$.


A zigzag we will say is a \textit{stickleback} iff its every zig and its every zag are tight. 

If $F=\zigzagTuple{f}{n}$ and 
$G=\zigzagTuple{g}{m}$ are zigzags such that the codomain of $f_n$
is the domain of $g_1$ so that the morphism $f_n \circ g_1$ is defined in \catcw then the zigzags $F$ and $G$ can be composed to give a zigzag 
$F::G=\tuple{f_1,f'_1,...f'_{n-1},f_n \circ g_1,g'_1, g_2,...g_m}$ 
of length $n+m-1$ as here:
$$
\arraycolsep=7pt
\left( \zigzag{w}{x}{f}{n}\right) :: \left(\zigzag{y}{z}{g}{m}\right)
$$

$$
=_{def} \  \zigzagfgcomposite
$$

\subsection{Comparison of Zigzags}

Say that zigzags are comparable if they are of the same length
and have identical object sequences.
If $F=\zigzagTuple{f}{n}$ and 
$G=\zigzagTuple{g}{n}$
 are comparable zigzags then define $F \leq G$
by $F \leq G$ iff for each $i$, $0 \leq i \leq n$
$f_i \leq g_i$ and for each $i$, $0 \leq i \leq n-1$ $f'_i \leq g'_i$.

\subsection{Tightening}

\newcommand{\tightlead}
  {f_1,f'_1,...f_{n-2},f'_{n-1},
       f_{n-1} \circ 
       \rng{
                \rst{f_n} \circ f'_{n-1}
                     }
}   

If $F$ is a zigzag, a \textit{tightening} of $F$ is a tight zigzag
$T \leq F$ such that for every tight zigzag $Z$ with $Z \leq F$,
we have $Z \leq T$.

Because $\leq$ is a partial order, any such tightening of $F$ is unique.

We will show by induction on length that every zigzag $F$ has a tightening and we will then write $T(F)$ for the tightening of $F$.
If the zigzag is of length $1$ then it has no zigs and no zags so is already tight. The inductive step is given by the next lemma: 

{  %lemma leftToRightTighteningSubLemma
   

\newcommand{\requiredTightening}
{\tuple{t_1,t'_1,...t'_{n-2},t_{n-1},\; 
    \rst{f_n} \circ f'_{n-1} \circ \rng{t_{n-1}},\; 
    \rst{f'_{n-1} \circ \rng{t_{n-1}}} \circ f_n
    }
}

\begin{lemma}
\llabel{leftToRightTighteningSubLemma}

Suppose  $F=\zigzagTuple{f}{n}$ is a zigzag in range category \catcw.
If  $\zigzagLeadingTuple{t}{n}$ is a tightening of the zigzag
\begin{equation}
\label{inductionPrecursor}
\tuple{\tightlead}
\end{equation}
then  zigzag
\begin{equation}
\label{requiredTightening}\requiredTightening
\end{equation}
is a tightening of $F$. 


\end{lemma}
\begin{proof}
We have to show first that (\ref{requiredTightening}) is tight. 

For this we need the following identities:
\begin{align}
\label{tightnessIdentityOne}
\rng{t_{n-1}} &= \rng{\rst{f_n} \circ f'_{n-1} \circ \rng{t_{n-1}}}  \\[0.1cm]
\label{tightnessIdentityTwo}
\rst{\rst{f_n} \circ f'_{n-1} \circ \rng{t_{n-1}}}&=\rst{\rst{f'_{n-1} \circ \rng{t_{n-1}}} \circ f_n}
\end{align}

To show that (\ref{tightnessIdentityOne}) holds 
we first note that by use of lemma \lref{fglemma} (iii) 
that its rhs simplifies to 
$\rng{\rst{f_n} \circ f'_{n-1}} \circ \rng{t_{n-1}}$ 
and that therefore  it
can be paraphrased as
\[
t_{n-1} \leq \rng{\rst{f_n} \circ f'_{n-1}}.
\]
This we show as follows
\newcommand{\stepOneReason}{
\parbox[t]{6cm}{
since
$t_{n-1} \leq  f_{n-1} \circ 
       \rng{
                \rst{f_n} \circ f'_{n-1}
                     }
$
as $t_{n-1}$ is the final element in the tightening
of (\ref{inductionPrecursor}),
} }
\newcommand{\stepThreeReason}{
\parbox[t]{6cm}{
because $\rng{f_{n-1}}$ is an idempotent
and 
for any idempotent $i$ and morphism $f$, $i \circ f \leq f$.
} }
\begin{align*}
\rng{t_{n-1}} &\leq  \rng{
       f_{n-1} \circ 
       \rng{
                \rst{f_n} \circ f'_{n-1}
                     }
                                      }
        && \mbox{\stepOneReason}\\
         &= \rng{f_{n-1}} \circ 
                      \rng{
                           \rst{f_n} \circ f'_{n-1}
                                    }
        && \mbox{by lemma \lref{fglemma} (iii),}\\[0.1cm]  
    &\leq  \rng{\rst{f_n} \circ f'_{n-1}}
        && \mbox{\stepThreeReason}
\end{align*}

Identity (\ref{tightnessIdentityTwo}) holds because both sides  simplify, using RR.3, and R.2, to $\rst{f'_{n-1} \circ \rng{t_{n-1}}} \circ \rst{f_n}$.

We can now show that (\ref{requiredTightening})
is the tightening of $F$ by showing that 
it dominates any tight zigzag $Z$ such that
$Z \leq F$ 
Assume such a $Z=\zigzagTuple{z}{n}$ and assume that

\[
\zigzagTuple{z}{n} \leq \zigzagTuple{f}{n}
\]
we need to show that 
\begin{equation}
\label{targetInequality}
Z \leq \requiredTightening.
\end{equation} 

First we note that  
\[
 z_{n-1} \leq f_{n-1} \circ 
       \rng{
                \rst{f_n} \circ f'_{n-1}
                     }
\]
because 
\begin{align*}
z_{n-1} &= z_{n-1} \circ \rng{\rst{z_n} \circ z'_{n-1}} && \mbox{trivially, since $Z$ is tight}\\
&\leq  f_{n-1} \circ 
       \rng{
                \rst{f_n} \circ f'_{n-1}
                     } &&\mbox{\parbox{3cm}{since $Z \leq F$.}}
\end{align*}

Thus since also $Z \leq F$ then we have that
\[
\zigzagLeadingTuple{z}{n} \leq \tuple{\tightlead}
\] 
and from this and the definition of $\zigzagLeadingTuple{t}{n}$
as a tightening of the rhs of this inequality we establish that
\[
\label{inequalityOne}
\zigzagLeadingTuple{z}{n} \leq \zigzagLeadingTuple{t}{n}
\]
 
To establish the target inequality (\ref{targetInequality})
we need 
in addition to show that
 \begin{equation}
 \label{inequalityTwo}
 z'_{n-1} \leq \rst{f_n} \circ f'_{n-1} \circ \rng{t_{n-1}}
 \end{equation}
 and
  \begin{equation}
 \label{inequalityThree}
 z_n \leq \rst{f'_{n-1} \circ \rng{t_{n-1}}} \circ f_n
 \end{equation}

The first of these holds because
\begin{align*}
z'_{n-1} &= \rst{z'_{n-1}} \circ z'_{n-1}  \circ \rng{z'_{n-1}} 
                && \mbox{trivially, by R.1 and RR.2,}\\
         &= \rst{z_n} \circ z'_{n-1}  \circ \rng{z_{n-1}} 
         &&\mbox{because $Z$ is tight,}\\
         &\leq \rst{f_n} \circ f'_{n-1}  \circ \rng{z_{n-1}}
         &&\mbox{because $Z \leq F$ ,}\\
         &\leq \rst{f_n} \circ f'_{n-1}  \circ \rng{t_{n-1}}
         &&\mbox{because $z_{n-1} \leq t_{n-1}$ is established
                               above in (\ref{inequalityOne}).}
\end{align*}

The second  because
\begin{align*}
 z_n &= \rst{z_n} \circ z_n
         &&\mbox{by R.1,}\\        
     &=  \rst{z'_{n-1}} \circ z_n 
         &&\mbox{since $Z$ is tight,}\\
     &\leq  \rst{\rst{f_n} \circ\ f'_{n-1} \circ \rng{t_{n-1}}} \circ f_n
         &&\mbox{by (\ref{inequalityTwo}) and since $z_n \leq f_n$,} \\
         &= \rst{f'_{n-1} \circ \rng{t_{n-1}}} \circ f_n
         && \mbox{by R.3, R.2 and R.1.}
\end{align*}
\end{proof}
} %end lemma leftToRightTighteningSubLemma
We have established that every zigzag $F$ has a unique tightening. 
We shall write $T(F)$ for the tightening of $F$. 


\begin{worktt}
In this lemma just given, if we denote by $F^L$
the zigzag $\tuple{\tightlead}$
and if we denote by $F^{L_r}$ the zigzag $\tuple{\rng{f'_{n-1}},f'_{n-1},f_n}$
then we can extract from  this lemma the following
\begin{observation}
\llabel{leftToRightTighteningObservation}
Suppose  $F=\zigzagTuple{f}{n}$ is a zigzag of length greater that $1$ in a range category \catcw
then there is a zigzag $F^L$ of length $n-1$ so that
$\tighten{F}= \tighten{F^L} \cnc \tighten{\rng{\tighten{F^L}}::F^{L_r}}$. 
\end{observation}
\end{worktt}

We also need the following takeaways from this lemma:
\begin{observation}
\llabel{bareTailObservation}
If $F=\zigzagTuple{f}{n}$ is a zigzag then
\[
T(F) = T(\zigzagBareTailTuple{f}{n}) ::\tuple{f_n}
\] 
\end{observation}
\noindent and from this there also follows:
\begin{observation}
\llabel{tailComposeObservation}
If  $F=\zigzagTuple{f}{n}$ is a zigzag 
and if the final morphism $f_n$ of $F$ composes with some morphism $g$
 so that the composite zigzag $F :: \tuple{g}$ is defined then
if $T(F)=\zigzagTuple{t}{n}$
and  $T(F :: \tuple{g})=\zigzagTuple{s}{n}$ then
\[
s_n = t_n \circ g
\]
\end{observation}

\begin{oldtt}
\workt{why?}
From this there follows:
\begin{observation}
\llabel{tightTailComposeObservation}
If  $F=\zigzagTuple{f}{n}$ is a tight zigzag 
and if the final morphism $f_n$ of $F$ composes with some morphism $g$
 so that the composite zigzag $F :: \tuple{g}$ is defined then
\[
\rng{T(F :: \tuple{g})}=\rng{f_n \circ g}
\]

\begin{worktt}
This changes in front of my eyes because where we could  end up if we want to go there is that
every zigzag has a restriction idempotent and a range idempotent defined by
\begin{align*}
\rst{Z} = \rst{\tighten{Z}}
\rng{Z} = \rng{\tighten{Z}}.
\end{align*} 
With such a notation then
\[
\rng{F :: \tuple{g}}=\rng{f_n \circ g}
\]
\end{worktt}
\end{observation}
\end{oldtt}
Lemma \lref{leftToRightTighteningSubLemma} describes how a zigzag can be tighten from left to right. We can also describe how a zigzag can be
tightened from right to left. This is what the next lemma is about.
{  %lemma rightToLeftTighteningSubLemma
  \newcommand{\tightTail}
  {\rst{f'_1 \circ \rng{f_1}} \circ f_2,\;
  f'_2,...f_{n-2},f'_{n-1}, f_n
  }      

\newcommand{\requiredRLTightening}
{\tuple{f_1 \circ \rng{\rst{t_2} \circ f'_1},\;
        \rst{t_2} \circ f'_1 \circ \rng{f_1},\; 
        t_2    ,...t_{n-1},t'_{n-1},t_n
    }
}
\begin{lemma}
\llabel{rightToLeftTighteningSubLemma}
Suppose  $F=\zigzagTuple{f}{n}$ is a zigzag in range category \catcw.
If  $\zigzagTrailingTuple{t}{n}$ is a tightening of the zigzag
\begin{equation}
\label{inductionRLprecursor}
\tuple{\tightTail}
\end{equation}
then  zigzag
\begin{equation}
\label{requiredRLTightening}\requiredRLTightening
\end{equation}
is a tightening of $F$. 
\end{lemma}
\begin{proof}
First of all we have to show that (\ref{requiredRLTightening}) is tight. 

For this we need the following identities:
\begin{align}
\label{RLtightnessIdentityOne}
\rng{f_1 \circ \rng{\rst{t_2} \circ f'_1}}&=\rng{\rst{t_2} \circ f'_1 \circ \rng{f_1}}   \\[0.1cm]
\label{RLtightnessIdentityTwo}
\rst{\rst{t_2} \circ f'_1 \circ \rng{f_1}}&=\rst{t_2}
\end{align}
The first of these follows from lemma \lref{fglemma} (i) and (iii).
The second because the lhs simplifies to 
$\rst{t_2} \circ \rst{f'_1 \circ \rng{f_1}}$
and so the identity can be paraphrased as 
$\rst{t_2} \leq \rst{f'_1 \circ \rng{f_1}}$ 
which we can prove as follows 

\newcommand{\stepOneReason}{
\parbox[t]{6cm}{because $t_2 \leq \rst{f'_1 \circ \rng{f_1}} \circ f_2$
as $t_2$ is the first element in the tightening of (\ref{inductionRLprecursor}),
}
}
\newcommand{\stepThreeReason}{
\parbox[t]{6cm}{because for any morphism $f$ 
                that can be composed with any idempotent $j$,
                    $f \circ j \leq f$.
}
}
\begin{align*}
\rst{t_2} &\leq \rst{\rst{f'_1 \circ \rng{f_1}} \circ f_2}
                && \mbox{\stepOneReason} \\
          &= \rst{f'_1 \circ \rng{f_1}} \circ \rst{f_2} 
                && \mbox{by R.3,} \\
          &\leq \rst{f'_1 \circ \rng{f_1}}
                && \mbox{\stepThreeReason}
\end{align*}


We can now show that (\ref{requiredRLTightening})
is the tightening of $F$ by showing that 
it dominates any tight zigzag $Z$ such that
$Z \leq F$ 
Assume such a $Z=\zigzagTuple{z}{n}$ and assume that

\[
\zigzagTuple{z}{n} \leq \zigzagTuple{f}{n}
\]
we need to show that 
\begin{equation}
\label{targetRLInequality}
Z \leq \requiredRLTightening.
\end{equation}

First we note that  
\[
 z_2 \leq \rst{f'_1 \circ \rng{f_1}} \circ f_2
\]
because 
\begin{align*}
z_2 &= \rst{z'_1 \circ \rng{z_1}} \circ z_2 
          && \mbox{trivially, since $Z$ is tight,}\\
    &\leq \rst{f'_1 \circ \rng{f_1}} \circ f_2 
     &&\mbox{\parbox{3cm}{since $Z \leq F$.}}
\end{align*}

Thus since also $Z \leq F$ then we have that
\[
\zigzagTrailingTuple{z}{n} \leq \tuple{\tightTail}
\] 
and from this and the definition of $\zigzagTrailingTuple{t}{n}$???
as a tightening of the rhs of this inequality we establish that
\[
\label{RLinequalityOne}
\zigzagTrailingTuple{z}{n} \leq \zigzagTrailingTuple{t}{n} ???
\]

To establish the target inequality (\ref{targetRLInequality})
we need 
in addition to show that
 \begin{equation}
 \label{RLinequalityTwo}
 z'_1 \leq  \rst{t_2} \circ f'_1 \circ \rng{f_1}
 \end{equation}
 and
  \begin{equation}
 \label{RLinequalityThree}
 z_1 \leq f_1 \circ \rng{\rst{t_2} \circ f'_1}
 \end{equation}
The first of these holds because
\begin{align*}
z'_1 &= \rst{z'_1} \circ z'_1  \circ \rng{z'_1} 
                && \mbox{trivially, by R.1 and RR.2,}\\
         &= \rst{z_2} \circ z'_1  \circ \rng{z_1} 
         &&\mbox{because $Z$ is tight,}\\
         &\leq \rst{z_2} \circ f'_1  \circ \rng{f_1}
         &&\mbox{because $Z \leq F$ ,}\\
         &\leq \rst{t_2} \circ f'_1  \circ \rng{f_1}
         &&\mbox{because $z_2 \leq t_2$ is established
                               above in (\ref{RLinequalityOne}).}
\end{align*}

The second because
\begin{align*}
z_1 &= z_1 \circ \rng{z_1}
        && \mbox{by RR.2,} \\ 
    &= z_1 \circ \rng{z'_1}
        && \mbox{since $Z$ is tight,} \\ 
    &\leq f_1 \circ \rng{\rst{t_2} \circ f'_1 \circ \rng{f_1}}
        && \mbox{by (\ref{RLinequalityTwo}) and since $z_1 \leq f_1$,} \\ 
    &= f_1 \circ \rng{\rst{t_2} \circ f'_1}
        && \mbox{by lemma \ref{fglemma} (iii), R.2, RR.1  and RR.2.}
\end{align*}
\end{proof}
}  % end of lemma rightToLeftTighteningSubLemma

From this lemma we cam make the following observations
\begin{observation}
\llabel{bareHeadObservation}
If $F=\zigzagTuple{f}{n}$ is a zigzag then
\[
T(F) =  \tuple{f_1} :: T(\zigzagBareHeadTuple{f}{n}) 
\] 
\end{observation}


\begin{observation}
\llabel{headComposeObservation}
If  $G=\zigzagTuple{g}{m}$ is a zigzag 
and if some morphism $f$ composes with the  initial morphism $g_1$ of $G$
 so that the composite zigzag $\tuple{f} :: G$ is defined then
if $T(F)=\zigzagTuple{t}{n}$
and  $T(\tuple{f} :: G)=\zigzagTuple{s}{n}$ then
\[
s_1 \leq f \circ t_1
\]
\end{observation}

\begin{observation}
\llabel{biggeristighter}
If we have two zigzags $F$ and $G$ and one is an extension of the other
then the tightening of the longer one truncated to be comparable with the smaller is less than or equal the tightening of the smaller one.\commentary{If needed should move earlier.}
\end{observation}

\begin{lemma}
\llabel{tighteningCommutesWithComposition}
If $F$ and $G$ are zigzags for which the composition $F::G$ is defined
then 
$ T(F :: G) = T(T(F) :: T(G)) $
\end{lemma}

\newcommand{\tightenedComposite}[1]
{
\zigzagOneAndAHalf{w}{x}{#1}
\begin{array}{c}
\dots \\[0.8cm]
\dots
\end{array}
\begin{array}{c c}
\Ob{x}{n} \\[0.8cm]
& \Ob{w}{n}=
\end{array}
\begin{arrows}
\ncarr{xn}{wn} \blabel{#1'_{n-1}}[0.2][0.4] %0.2 was 0.45
\end{arrows}
\kern-15pt
\arraycolsep=7pt
\zigzagFlex{y}{z}{#1}{n}{n+m}
}

\begin{proof}
Suppose
\begin{itemize}
\item 
$F=\zigzagTuple{f}{n}$ is a zigzag shaped like this
$$
\arraycolsep=7pt
\zigzag{w}{x}{f}{n}
$$
\item $G=\zigzagTuple{g}{m}$ is a zigzag shaped like this
$$
\arraycolsep=7pt
\zigzag{y}{z}{g}{m}
$$
such that $x_n = y_1$ so that $F::G$ is defined,
\item $T(F)=\zigzagTuple{t}{n}$,
\item $T(G)=\zigzagTuple{s}{m}$,
\item $T(T(F) :: T(G))=\zigzagTuple{r}{n+m}$ so that we have
$$
\arraycolsep=7pt
\tightenedComposite{r},
$$
\item Let $Q$ be the tightening of $F :: G$ 
and suppose $Q$ to be the zigzag $\zigzagTuple{q}{n+m}$ 
so that we have
$$
\arraycolsep=7pt
\tightenedComposite{q}.
$$
\end{itemize}

We will prove the result by showing 
\begin{enumerate}[(a)]
\item that $T(T(F) :: T(G)) \leq T(F :: G)$ and
\item that $T(F :: G) \leq T(T(F) :: T(G))$.
\end{enumerate}

To establish (a), since, by definition, $T(T(F) :: T(G))$ is tight
it suffices to show that 
$T(T(F) :: T(G)) \leq F:: G$. 
Since $T(T(F) :: T(G)) \leq T(F) :: T(G)$, it therefore suffices to show that
\[
T(F) :: T(G) \leq F:: G
\]
i.e. that 
\[
\tuple{t_1,t'_1,...t'_{n-1},s_1 \circ t_n,s'_1,...s_m}
\leq \tuple{f_1,f'_1,...f'_{n-1},f_n \circ g_1 ,g'_1,...g_m}
\]
This follows because
\[
\zigzagTuple{t}{n} \leq \zigzagTuple{f}{n},
\]
as the left zigzag is defined to be a tightening of the right, and similarly
\[
\zigzagTuple{s}{m} \leq \zigzagTuple{g}{m}.
\]
Moreover, since $t_n \leq f_n$ and $s_1 \leq g_1$, it follows by Lemma~\lref{orderingsCompose} that
\[
t_n \circ s_1 \leq f_n \circ g_1.
\]

To establish (b), because $T(F :: G)$ is tight it suffices to show that
$T(F :: G) \leq T(F) :: T(G)$. In other words that
\begin{equation}
\label{pointwiseTargetInequality}
\tuple{q_1,q'_1,...q'_{n-1},q_n,q'_n,...q_{n+m}}\leq \tuple{t_1,t'_1,...t'_{n-1},t_n \circ s_1,s'_1,...s_m}
\end{equation}

We prove this piecemeal
starting with the those inequalities  
that precede the inequality for $q_n$.
First note that $\rst{q_n} \leq \rst{f_n}$
because as part of the tightening of $F \circ G$ we have that
$q_n \leq f_n \circ g_1$ and therefore 
$\rst{q_n} \leq \rst{f_n \circ g_1} \leq \rst{f_n}$. 

From this, and again because of the definition of $Q$ as the tightening
of $F \circ G$, we have that 
\[
\zigzagBareTailTuple{q}{n} \leq \zigzagBareTailTuple{f}{n}
\]
from which follows 
\[
\zigzagBareTailTuple{q}{n}  \leq T(\zigzagBareTailTuple{f}{n})
\]
and therefore
\begin{align*}
\zigzagBareTailTuple{q}{n} \circ f_n  &\leq T(\zigzagBareTailTuple{f}{n}) \circ f_n,\\
                                 &=  \tuple{t_1,... t_n} 
                                 \mbox{ by observation \ref{bareTailObservation},} \\
\end{align*}
which gives us the required inequalities for those $q_i$ and $q'_i$ up to and including $q'_{n-1}$.

Similar reasoning, using observation \ref{bareHeadObservation}, can be used to establish
the inequalities for all those $q_i$ and $q'_i$ that follow after the inequality for
$q_n$. The remaining inequality to be established is that
\[
\label{qleqtnsOne}
q_n \leq t_n \circ s_1.
\]
To show that this holds first consider the zigzag $F :: \tuple{g_1}$.
This is the initial segment of the zigzag $F :: G$. 
Because $\tuple{q_1,...q_{n+m}}$ is the tightening of 
$F :: G$ it follows
using observation \lref{biggeristighter}  that 
\begin{equation}
\label{qTFg}
\tuple{q_1,...q'_{n-1},q_n} \leq T(F :: \tuple{g_1}).
\end{equation} 
By observation \lref{tailComposeObservation} 
the final element of $T(F :: \tuple{g_1})$ is $t_n \circ g_1$
and therefore we have in particular that
\begin{equation}
\label{qgtn}
q_n \leq t_n \circ g_1.
\end{equation}
By similar considerations and using observation \lref{headComposeObservation}
we have that 
\begin{equation}
\label{qTfG}
\tuple{q_n, q'_n,...q_{n+m}} \leq T(\tuple{f_n} :: G)
\end{equation}
and 
\begin{equation}
\label{qsfn}
q_n \leq f_n \circ s_1.
\end{equation}
Finally, using lemma \lref{pointwiseorderingsCompose} we establish  
(\ref{qleqtnsOne}) from
(\ref{qgtn}) and (\ref{qsfn}), as required.
\end{proof}


\subsection{The category of tight zigzags}
\llabel{categorytightzigzags}
If \catcw is a range category then we can define the category
$Z_C$ to be the category whose objects are the objects of \catcw
and such that $Hom_{Z_C}(a,b)$ is the set of tight zigzags of any length $n$
$$
\arraycolsep=7pt
\zigzag{x}{y}{f}{n}
$$
such that $x_1=a$ and $y_n=b$.


Composition in $Z_C$ is defined by composing the zigzags and then tightening  thus
\[
F \circ G =_{def} T(F :: G).
\]
Composition is associative since 
\begin{align*}
(F \circ G) \circ H 
                &= T(T(F \conc G) \conc H) &&   \mbox{by definition,} \\
                &= T(T(F \conc G) \conc T(H))
                        &\mbox{since $H$ is tight,}\\
                &= T((F \conc G) \conc H)  
                        &\mbox{by lemma \lref{tighteningCommutesWithComposition},}\\
                &= T(F \conc (G \conc H))
                        &\mbox {since $\conc$ is associative,} \\
                &= T(T(F) \conc T(G \conc H))       
                        &\mbox{by lemma \lref{tighteningCommutesWithComposition},}\\
                &= T(F \conc T(G \conc H))))
                        &\mbox{since $F$ is tight,} \\
                &= F \circ (G \circ H) &&   \mbox{by definition.}                 
\end{align*}

Identity morphisms in $Z_C$ consist of the identity morphisms of \catcw considered 
as zigzags on length 1.

To show that $Z_C$ is a category with support we need the following lemma:
\begin{lemma}
\llabel{midpointLemma}
If $F=\zigzagTuple{f}{n}$ and $G=\zigzagTuple{g}{m}$ are tight zigzags
then 
\[
F \circ G = (F \circ \tuple{\rst{g_1}}) :: (\tuple{\rng{f_n}} \circ G)
\]
\end{lemma}
\begin{proof}
\begin{align*}
F \circ G &= T(\zigzagTuple{f}{n}::\zigzagTuple{g}{m})
           &&\mbox{definition of $f \circ G$,}\\
          &= T(\zigzagTupleLedOut{f}{n}{\rst{g_1}} :: \zigzagTupleLedIn{\rng{f_n}}{g}{m} )
          && \mbox{since $f_n \circ g_1 = f_n \circ \rst{g_1} \circ \rng{f_n} \circ g_1$,} \\
          &=  T(T(\zigzagTupleLedOut{f}{n}{\rst{g_1}}) :: 
                           T(\zigzagTupleLedIn{\rng{f_n}}{g}{m}) )
          && \mbox{by lemma \lref{tighteningCommutesWithComposition},}\\
          &= T(\zigzagTupleLedOut{f}{n}{\rst{g_1}}) :: 
                           T(\zigzagTupleLedIn{\rng{f_n}}{g}{m}) 
          && \mbox{by lemma was \lref{fglemma}(iv)
                        now lemma \lref{fgJoinLemma},}\\
          &= (F \circ \tuple{\rst{g_1}}) :: (\tuple{\rng{f_n}} \circ G)
          && \mbox{by definition of $\circ$.}
\end{align*}
\end{proof}

\begin{lemma}
\llabel{howToTightenAConcatenation}
If $F$ and $G$ are zigzags that compose so that
$F \cnc G$ is defined then
\begin{align*}
\tighten{F \cnc G} 
    &= 
    \Tighten{ \tighten{F} \cnc \rst{\tighten{G}}}
     \cnc
    \Tighten{ \rng{\tighten{F}} \cnc \tighten{G}} \\
    &= 
    \Tighten{ F \cnc \rst{\tighten{G}}}
     \cnc
    \Tighten{ \rng{\tighten{F}} \cnc G} \\
\end{align*}
\end{lemma}
\begin{proof}
See daybook 23rd June 2026 for a partial proof that simply uses definition of tightening. Can probably use midpointLemma also but this proof should be neater and maybe midpointLemma follows from this lemma.
\end{proof}
\newcommand{\abrst}[1]{\rst{\tighten{#1}}}
\newcommand{\abrng}[1]{\rng{\tighten{#1}}}

\begin{lemma}\commentary{where should these be in the pecking order?}
\commentary{needed before associativity of }
\llabel{FGsupportLemmas}
If $F$ and $G$ are zigzags such that $F \cnc G$ is defined then
\commentary{by $\rst{ZZ}$ I am inventing terminology on the fly. I mean the singleton zigzag $\tuple{\rst{zz_1}}$ where
$zz=\tuple{zz_1,zz'_1,...z_n}$ for some $n$.
}
\commentary{which of these do I need?}
\begin{enumerate}[(i)]
\item \llabel{leftpushin} 
$\tighten{F \cnc G} = \tighten{F \cnc \abrng{F} \cnc G}$
\item \llabel{rightpushin}
$\tighten{F \cnc G} = \tighten{F \cnc \abrst{G} \cnc G}$
\item \llabel{bothpushin}
$\tighten{F \cnc G} = \tighten{F \cnc \abrng{F} \cnc \abrst{G} \cnc G}$
\item \llabel{abrstpushin}
$\abrst{F \conc G} = \abrst{F \conc \abrst{G}}$
\item \llabel{arbrngpushin}
$\abrng{F \conc G} = \abrng{\abrng{F} \conc G}$
\end{enumerate}
\end{lemma}

\begin{lemma}[howToTightenAThreefoldConcatenation]
\llabel{howToTightenAThreefoldConcatenation}
If $F$, $Z$ and $G$ are zigzags that compose so that
$F \cnc Z \cnc G$ is defined then
\[
\Tighten{F \cnc Z \cnc G} = 
    \tighten{ F \cnc \abrst{{Z \cnc \abrst{G}}}}
     \cnc 
     ZZ
     \cnc
    \tighten{ \abrng{F \cnc Z} \cnc G}
\]
where
\[
ZZ=\Tighten{\abrng{F}
              \cnc
            Z
              \cnc
            \abrst{G}
           }
\]
We can use \lref{FGsupportLemmas} \lref{abrstpushin} or \lref{abrngpushin}
to make this symmetric but which of the two possible symmetric statements do I want?
\end{lemma}
\begin{proof}

Hand written on loose A4. 
First note that
\begin{align*}
\tighten{ F \cnc Z \cnc G} 
               &=\tighten{ (F \cnc Z) \cnc G} \\ 
               &= \tighten{ (F \cnc Z) \cnc \abrst{G}} 
                     \cnc
                  \tighten{ \abrng{F \cnc Z} \cnc G} 
               &&\mbox{by lemma 
                    \lref{howToTightenAConcatenation}.}
\end{align*}
Then that 
\begin{align*}
\tighten{ (F \cnc Z) \cnc \abrst{G}}  
        &= \tighten{ F \cnc (Z \cnc \abrst{G})} \\
        &= \tighten{ F \cnc \abrst{{Z \cnc \abrst{G}}}} 
              \cnc
           \tighten{ \abrng{F} \cnc (Z \cnc \abrst{G})} 
           &&\mbox{by lemma 
                    \lref{howToTightenAConcatenation}.} \\
        &= \tighten{ F \cnc \abrst{{Z \cnc \abrst{G}}}} 
              \cnc
           ZZ 
           &&\mbox{where we define $ZZ$ as above.} 
\end{align*}   
\end{proof}

\begin{lemma}
\llabel{howToTightenAThreefoldConcatenationTwo}
If $F$, $Z$ and $G$ are zigzags that compose so that
$F \cnc Z \cnc G$ is defined then
\commentary{\highlight{review why} I want this}
\commentary{\highlight{use this next!}}
\[
\Tighten{F \cnc Z \cnc G} = 
    \Tighten{ F \cnc \rst{ZZ}}
     \cnc 
     ZZ
     \cnc
    \Tighten{ \rng{ZZ} \cnc G}
\]
where
\[
ZZ=\Tighten{\rng{\tighten{F}}
              \cnc
            Z
              \cnc
            \rng{\tighten{G}}
           }
\]
\end{lemma}
\begin{proof}
This restatement of previous lemma follows from
\lref{FGsupportlemmas} by 
use one or two of \lref{leftpushin}, \lref{rightpushin}, \lref{bothpushin}.
\end{proof}
\subsection{The category of tight zigzags $Z_C$ is a category with support}

If $F=\zigzagTuple{f}{n}$ is a morphism in $Z_C$ then define $\rst{F}$ to be $\tuple{\rst{f_1}}$.

We need show

R.1 For $F:a \morph b$ in $Z_C$ $$\rst{F} \circ F =F$$.

R.2. If \FGsourcediag in $Z_C$ then
$$\rst{G} \circ \rst{F}=\rst{F} \circ \rst{G}.$$

R.3. If \FGsourcediag in $Z_C$ then
$$\rst{\rst{F} \circ G} = \rst{F} \circ \rst{G}$$.

wR.4. If $\sequentialdiag{a}{b}{c}{F}{G}$ in $Z_C$ then
\[
\rst{F \circ \rst{G}} = \rst{F \circ G}.
\]

R.1 follows because in \catcw, $f_0 \circ \rng{f_0} = f_0$. 

R.2 follows because in \catcw $\rng{f_0} \circ \rng{g_0} = \rng{g_0} \circ \rng{f_0}$.

R.3 follows because
\commentary{work on this proof}
\begin{align*}
\rst{\rst{F} \circ G} 
                   &= \rst{T(\tuple{g_0 \circ \rng{f_0},g'_0,...g_m})}\\
                   &= \rst{\tuple{g_0 \circ \rng{f_0},s'_0,...s_m}}
                                             &&\mbox{\parbox{5cm}{for some $s'_0,...s_m$
                                                       by some new lemma about tightening when only head position loose,}}\\
                   &=\tuple{\rng{g_0 \circ \rng{f_0}}} 
                                            && \mbox{by definition}\\
                   &=\tuple{\rng{g_0} \circ \rng{f_0}}  
                                            && \mbox{by lemma \lref{fglemma} (iii)}\\
                   &=\rst{F} \circ \rst{G}      
\end{align*}

If $F$ is of length 1 then wR.4 holds because
it holds in the underlying range category \catcw. 
If $F$ is of length $n > 1$ then wR.4 follows because lemma
\lref{midpointLemma} informs us that $\rst{f \circ G}$ depends only on the $\rst{g_1}$ component of $G$ i.e. only on $\rst{G}$.


\subsection{Shortening and Short Sticklebacks}

If $F=\zigzagTuple{f}{n}$ is a zigzag and 
if for some $i$,
$1 \leq i \leq n-1$,
there is a morphism
$s$ such that $f'_i \circ s = f_{i+1}$ so that the triangle in 
\[
\zigzagShortenable{x}{y}{f}{n}{i}{s}
\]
commutes then the zigzag 
\[
\arraycolsep=7pt
\zigzagShortened{x}{y}{f}{n}{i}{s}
\]
is said to be a \textit{shortening} of $F$ and is said to be \textit{at position i of F}.

\newcommand{\rngTightener}[2]
    {
          \rng{
                \rst{#1_{#2+1}} \circ #1'_{#2}
               }            
    }

\newcommand{\rngTightenerOne}[1]
    {
          \rng{
                \rst{#1_{2}} \circ #1'_{1}
               }            
    }

\newcommand{ \rstTightener}[2]
    {
          \rst{
                #1'_{#2} \circ \rng{#1_{#2}}  
               }
    } 
\begin{notebox}[Change to the above June 24 2026]
Change the above. If $F$ is \emph{not} tight
then need the shortening to be not $f_i \circ s$ as is shown in the diagram but to be \oldt{$f_i \circ \rngTightener{f}{i} \circ s$}  
\newt{$f_i \circ \rng{f'_i} \circ s$}.
\end{notebox}

Significantly, the shortening of a tight zigzag is  tight.
It is easy to check this.

Cockett shows that if the category \catcw is a range category
satisfying the addition condition RR.5 then every tight
zigzag of morphisms in \catcw has a unique shortening. 

A stickleback which cannot be shortened is said to be
a \textit{short stickleback}.

\begin{lemma}
\llabel{shorteningCharacterisationLemma}
if $Z$ is a zigzag of length $2$ and $s$ is a shortening
of $Z$ as shown here
\[
\zigzagTwo{x}{y}{z}
\begin{arrows}
\ncarr[15]{y1}{y2} \alabel{s}
\end{arrows}
\]
then
and if $\tuple{zz}$ is the shortening of $Z$
so that
$zz=z_1 \circ \rng{z'_1} \circ s$
then
\begin{enumerate}[(i)]
\item $\rst{zz}= \rst{z_1 \circ \rngTightenerOne{z}}$,
\item $\rng{zz} =\rng{\rstTightener{z}{1} \circ z_2  }$.
\end{enumerate}
\end{lemma}
\begin{proof}
\begin{enumerate}[(i)]
\item Suffices to show that 
$\rst{\rng{z'_1} \circ s}=\rst{\rngTightenerOne{z}}$
which we do as follows
\begin{align*}
\rst{\rng{z'_1} \circ s} 
    &= \rst{\rng{z'_1} \circ \rst{s}} && \mbox{wR.4}\\
    &= \rst{ \rng{ z'_1 \circ \rst{s} } }
                                      && \mbox{RR.3} \\
    &= \rst{ \rng{ \rst{z'_1 \circ s} \circ z'_1 } }
                                      && \mbox{R.4} \\
        &= \rst{ \rng{ \rst{z_2} \circ z'_1 } }
                        && \parbox[t]{5cm}{from the initial assumption 
                        $z'_1 \circ s = z_2$.} 
\end{align*}
Now using wR.4 we get
$\rst{zz}
=\rst{z_1 \circ \rng{z'_1} \circ s}
=\rst{z_1 \circ \rst{\rng{z'_1} \circ s}}
=\rst{z_1 \circ \rst{\rngTightenerOne{z}}}
= \rst{z_1 \circ \rngTightenerOne{z}}$.
\item \begin{align*}
\rng{zz} &= \rng{z_1 \circ \rng{z'_1} \circ s}\\
         &= \rng{
                    \rstTightener{z}{1}
                    \circ z'_1 \circ s
                 }
            &&\mbox{by R.4 and since, by RR.1,
            $ \rng{z'_1}=\rst{\rng{z'_1}}$,
            }\\
         &= \rng{\rstTightener{z}{1} \circ z_2 }
            &&\parbox[t]{5cm}{from the initial assumption  $z'_1 \circ s = z_2$.} 
\end{align*}
\end{enumerate}
\end{proof}

\begin{lemma}
\llabel{secondShorteningCharacterisationLemma}
\end{lemma}
\begin{proof}
If $Z$ is a zigzag of length $2$ and $s$ is a shortening of $Z$
then
\begin{enumerate}[(i)]
\item 
$\absrst{\shorten{i}{s}{Z}} = \abrst{Z}$
\item
$\absrng{\shorten{i}{s}{Z}} = \abrng{Z}$
\end{enumerate}
\end{proof}

\subsection{Prelude to a category of short zigzags}

In what follows lemma 
\lref{deferredShorteningLemma}
describes the interaction of shortening and tightening 
and is the key to showing that there is a category of short sticklebacks. Before we get to this lemma we have a series of
subordinate lemmas that are required for its proof. 
 
\begin{lemma}
\llabel{shorteningSpecialisationLemma}
In any restriction category \catcw, 
if $f' \leq g'$ and $f \leq g$ then if  
$\nudgeup{1.5cm} % moves the triangle down to make room for arc
\downTriangleShape{x_i}{y_{i-1}}{y_i}
\downToLeft{g'}[0.6]
\downToRight{g}
\leftToRightArcing[a]{s}{30}{30} 
$ commutes 
and $\rst{f}=\rst{f'}$ then
$\downTriangleShape{x_i}{y_{i-1}}{y_i}
\downToLeft{f'}[0.6]
\downToRight{f}
\leftToRightArcing[a]{s}{30}{30} 
$ commutes.
\end{lemma}
\begin{proof}
\begin{align}
f' \circ s 
   &= \rst{f'} \circ g'\circ s 
       && \mbox{because $f' \leq g'$,}\\
   &= \rst{f} \circ g 
   && \mbox{since $\rst{f'}=\rst{f}$ 
              and $g' \circ s = g$, }\\
   &= f
   &&\mbox{since $f \leq g$}
\end{align}
\end{proof}

\begin{corollary}
\llabel{shorteningAtightenedRestriction}
If $F$ is a zigzag and $s$ is a shortening of $F$ at position $i$,
if $j$ is a restriction endomorphism 
such that $F :: \tuple{j}$ is defined then
$s$ is a shortening of 
the tightening of $F :: \tuple{j}$ at position$i$.
\end{corollary}

\begin{aside}
If I talk about shortening then the context
needs to be \highlight{RR.5} range category.
\end{aside}

\begin{aside}
I can write  $F= F^{L} \cnc F^{L_r}$

and thus $F\cnc\tuple{g} = (F\cnc\tuple{g})^{L} \cnc (F\cnc\tuple{g})^{L_r}$??

\end{aside}

From this lemma we can see that if a zigzag $F$ shortens then $T(F)$ shortens.
\newpage
\begin{lemma}
\llabel{associativityLeftSubLemma}

If $F$ is a tight zigzag of length $n$ and if $s$ is a shortening of $F$ at 
position $i$ where
$i \leq  n-2$ and if $g$ is a morphism such that $F :: \tuple{g}$ is
defined then $s$ is a shortening of $\tighten{F\cnc\tuple{g}}$
and 
\[
\tighten{
  \shorten{F\cnc\tuple{g}}
          {s}{i}
        }
=\Shorten{\tighten{F\cnc\tuple{g}}}
         {s}{i}
\]
\end{lemma}
\begin{worktt}
\begin{proof}
\newcommand{\Fg}{F \cnc \tuple{g}}
\newcommand{\FgL}{(\Fg)^L}
\newcommand{\FgLr}{(\Fg)^{L_r}}
\newcommand{\shortendFg}{\shorten{F\cnc\tuple{g}}{s}{i}}
\newcommand{\shortendFgL}{\bigl(\shortendFg\bigr)^L}
\newcommand{\shortendFgLr}{\bigl(\shortendFg\bigr)^{L_r}}
\newcommand{\FgLshortened}{\shorten{\FgL}{s}{i}}
\newcommand{\FgLtightened}{\tighten{\FgL}}
\newcommand{\Z}{\Tighten{\shortendFgL}}
\newcommand{\Zrewritten}{\shorten{\FgLtightened}{s}{i}}
\newcommand{\ZZ}{\FgLtightened \cnc \tighten{
                     \rng{Z} \cnc \FgLr
                        }
}
\newcommand{\ZZrewritten}{\FgLtightened \cnc \tighten{
                     \rng{\FgLtightened} \cnc \FgLr
                        }
}
First note that since $i \leq n-2$ yt follows that
$s$ is a shorting of $F\cnc \tuple{g}$ at position $i$. From this it follows by corollary 
\lref{shorteningAtightenedRestriction} that $s$ is also a shortening 
at  position $i$ of $\tighten(F\cnc \tuple{g})$. 

Now the left hand side of the required identity can be rewritten using 
observation \ref{leftToRightTighteningObservation} as
\begin{align*}
\text{LHS} 
    &= Z \cnc \tighten{
                     \rng{Z} \cnc \shortendFgLr
                        }
    &          &&\parbox{6cm}{where $Z=\Z$,}
\end{align*}
\begin{newtt}

\end{newtt}
\begin{oldtt}
but to take the front end of the non-final 
shortening of a zigzag is to shorten the front end i.e. 
\[
\shortendFgL = \FgLshortened 
\] \commentary{the is first identity isn't true when $i = n-2$}
and shortening the front end does not change the final elements of the zigzag therefore
\[
\shortendFgLr = \FgLr
\]
\end{oldtt}


Therefore we can rephrase $Z$ and use the inductive hypothisis to get
\begin{align*}
\text{LHS} &= Z \cnc \tighten{
                     \hat{Z} \cnc \FgLr
                        }
              &&\parbox{6cm}{where $Z=\Zrewritten$,}\\
           &= \Shorten{\ZZ}{s}{i} && \parbox{6cm}{because shortening can moves outside the concatenation $\cnc$,}\\
           &= \Shorten{\ZZrewritten}{s}{i} 
                   && \parbox{6cm}{observation xxx that $\rng{\shorten{ZZ}{s}{i}}=\rng{ZZ}$,}\\
           &= \Shorten{\tighten{F\cnc\tuple{g}}}
                     {s}{i} 
            && \parbox{6cm}{by use of observation \ref{leftToRightTighteningObservation} again,}\\
         &= \text{RHS.}
\end{align*}
\end{proof}
\end{worktt}
\begin{lemma}
\llabel{associativityMidLemma}
\begin{newtt}
If $F$ is a tight zigzag of length $n$ and if $s$ is a shortening of $F$ at 
position $n-1$, if $g$ is a morphism such that $F :: \tuple{g}$ is
defined then $s \circ g$ is a shortening of $\tighten{F\cnc\tuple{g}}$
and 
\[
\tighten{
  \shorten{F\cnc\tuple{g}}
          {s}{i}
        }
=\Shorten{\tighten{F\cnc\tuple{g}}}
         {s}{i}
\]
\end{newtt}
\end{lemma}
\begin{lemma}
\llabel{associativityRightSubLemma}
\end{lemma}

\begin{lemma}
\llabel{associativitySublemma}
Suppose $F$ of length $n$ and $G$ of length $m$ are tight zigzags 
in an RR.5 range category \catcw 
and suppose that $s$ is a shortening
of $F::G$ at some index $i_s$, $1 \leq i_s \leq n +m - 1$, then
\begin{equation}
\label{associativitySublemmatarget}
T(shorten(F::G,i_s,s))=shorten(T(F::G),i_s,s))
\end{equation}
\end{lemma}
\begin{proof}
\begin{newtt}\commentary{new proof 24 May 2026}
This lemma follows directly from lemmas 
\lref{associativityLeftSubLemma},
\lref{associativityMidSubLemma} and
\lref{associativityRightSubLemma} and from
lemma \lref{midpointLemma}.

Note that the RHS of the target identity can be rwritten as
\[
shorten(F \circ G,i_s,s)
\]
Hmmm
Isn't $F::G$ equal
\[
(F::\rst{G})::(\rng{F} :: G)
\]
and doesn't this mean that $shorten(F::G,i_s,s)$ is
either
\[
shorten(F::\rst{G},i_s,s) ::(\rng{F}::G)
\]
or
\[
(F::\rst{G}) :: shorten(\rng{F}::G),i_s - n, s)
\]

\end{newtt}
\workt{proof was commented out when it went out of date}

\workt{This proof will be replaced by a coming together of two earlier lemmas.}

\workt{associativityLeftSubLemma}
\workt{associativityRightSubLemma}
\begin{aside}
19 May 2026.
The definition of tightening has been modified since this proof was writtem.
We have right to left and left to right versions. The proof should be revised accordingly.
WE no linger need the framedEquations environment so I will remove these
out of date equations.
\end{aside}
In such a situation, $F::G$ is a zigzag in which 
elements before the element $f_n \circ g$ need 
tightening to the left and terms to after need tightening to the right. 
I am going to refer to  the leading $2n -1$ elements of T(F::G)
as $ff_1, ff'_1,...ff'_{n-1}$ and $ff'_n$ and I am going to refer 
to the trailing $2m-1$ elements as $gg_1,gg'_1,...gg_m$.
What  this means is that the element that is refereed to as
$ff_n$ is also referred to as $gg_1$. 
With this way of referencing the elements of $T(F::G)$ 
then I can describe $T(F::G)$ as follows:
\tbd
 

There are three cases to consider.
\paragraph{Case (i) --- $i_s \leq  n-2$}
Use \lref{associativityLeftSubLemma}. Plus some.

\paragraph{Case (ii) --- $i_s = n-1$}
In this case the shortening of the lhs of (\ref{associativitySublemmatarget}) eliminates the pair
$f'_{n-1}, f_n \circ g$ 
based on the assumption that this diagram commutes:
\begin{equation}
\label{case.ii.lhsshortening} 
\begin{array}{c c c c}
           &           &           &\Ob{z}{1} \\[0.8cm]
\Obp{x}{n} &           & \Ob{x}{n}            \\[0.8cm]
           & \Ob{w}{n}. &  
\end{array}
\begin{arrows}
\drawMorphismArcing{xnp}{z1}{a}{s}{15}{15}[0.5]
\ncarr{xn}{z1}\blabel{g_1}
\ncarr{wn}{xnp}\alabel{f'_{n-1}}
\ncarr{wn}{xn}\blabel{f_n}
\end{arrows}
\end{equation}

The shortening of the rhs of (\ref{associativitySublemmatarget}) eliminates pair
$ff'_{n-1}, ff_n$
and is supported by this diagram
\begin{equation}
\label{case.ii.rhsshortening} 
\begin{array}{c c c c}
           &           &           &\Ob{z}{1} \\[0.8cm]
\Obp{x}{n} &           &                      \\[0.8cm]
           & \Ob{w}{n} &  
\end{array}
\begin{arrows}
\drawMorphismArcing{xnp}{z1}{a}{s}{15}{15}[0.5]
\ncarr{wn}{xnp}\alabel{ff'_{n-1}}
\ncarr{wn}{z1}\blabel{ff_n}
\end{arrows}
\end{equation}
whose commutivity follows 
by lemma \lref{shorteningSpecialisationLemma}.

The shortening  of $F::G$ by $s$ is 
$\tuple{f_1,f'_1,...f'_{n-1},
        f_{n-1} \circ s,
        g'_1, g_2,...g'_{m-1},g_m
       }$
The tightening of this shortening
is described by appropriately varying above
equations (\ref{tighteningcomposite:ffn}) to (\ref{tighteningcomposite:ggj}).
In the tightening of the shortening there is one less element
than in the original tightening. In the tightening the replacement element $f_{n-1} \circ s$ is unchanged by tightening and we refer to it as both  $tsf_{n-1}$ and $tsg_1$. There is one less leading elements and the same number of trailing elements. 
The tightening is as follows:
\tbd
On the other hand  the tightening, $T(F::G)$ 
is described by equations 
(\ref{tighteningcomposite:ffn}) -- 
(\ref{tighteningcomposite:ggj}), above and this
zigzag can be shorten by $s$ because, by lemma \ref{shorteningSpecialisationLemma},

\[
\downTriangleShape{w_n}{x_{n-1}}{z_1}
\downToLeft{ff'_{n-1}}[0.6]
\downToRight{ff_n=gg_1=f_n \circ g_1}
\leftToRightArcing[a]{s}{15}{15}
\]
commutes. 
The shortening of this tightening is the following variation
of equations (\ref{tighteningcomposite:ffn}) to (\ref{tighteningcomposite:ggj})

\tbd

To establish that $T(shorten(F::G,n,s)=shorten(T(F::G),n,s))$
we have to resolve the difference between equations 
(\ref{tighteningShortenedComposite:ffnp}) and
(\ref{shortenedTightenedComposite:stfnp}) i.e. we have to prove that $ff'_{n-1} \circ s
= f'_{n-1} \circ s$ which we can do as follows:
\begin{align*}
ff'_{n-1} \circ s 
   &=\rst{ff_n}\circ f'_{n-1}\circ s&&\mbox{using  (\ref{tighteningcomposite:ffpi}),}\\
   &=\rst{f_n\circ g_1}\circ f'_{n-1}\circ s&&\mbox{using (\ref{tighteningcomposite:ffn}),}\\
   &= \rst{f'_{n-1}\circ s}\circ f'_{n-1}\circ s&&\mbox{using commutivity of (\ref{case.ii.lhsshortening}),}\\
   &=f'_{n-1}\circ \rst{s} \circ s&&\mbox{by R.4,}\\
   &= f'_{n-1} \circ s     && \mbox{as required, by R.1.}
\end{align*}

\begin{aside}
just noticed in early parts of this proof need to change x to w and y to x
\end{aside}
\paragraph{Case (iii) --- $i_s \geq n$}
Use \lref{associativityRightSubLemma}. Plus some.
\end{proof}
\newpage

\begin{lemma}
\llabel{TSExtractionLemma}
If $F$ is a zigzag and $s$ is a shortening of $F$ at a position $i$ then
\[
\tighten{\shorten{i}{s}{F}} = \shorten{i}{s}{\tighen{F}}
\]
\end{lemma}
\begin{proof}
$F$ can be expressed as $F_L \cnc Z \cnc F_R$ where 
$Z$ is of length $2$ and where $s$ is a shortening of $Z$ at (its only)
position $1$. We can then use the tightening of $Z$ and lemma 

\end{proof}
\begin{aside}
The following lemma suggests that the normal form for a zigzag is the
shortening of the sightening.
\end{aside}
\begin{lemma}
\llabel{deferredShorteningLemma}
If $Z_1$, $Z_2$ are tight zigzags in
\newt{an RR.5 range category} such that the composition $Z_1 :: Z_2$ is defined then
\begin{enumerate}[(i)]
\item 
$shorten(tighten(shorten(Z_1) :: Z_2))= shorten(tighten(Z_1 :: Z_2)).$
\item
$shorten(tighten(Z_1 :: shorten(Z_2)))= shorten(tighten(Z_1 :: Z_2)).$
\end{enumerate}
\end{lemma}
\begin{proof}
Since in both (i) and (ii) the innermost shortening is a combination 
of individual shortenings and since
it follows by lemma \lref{associativitySublemma} that 
each of these individual shortenings can be moved outside of the tightening.
\end{proof}

\subsection{The category of short sticklebacks}

If \catcw is an RR.5 range catgeory then
 there is a category of short sticklebacks over \catc, which we denote
 $\Stk$.
Composition of short sticklebacks is defined by
\[
F \circ G = shorten(T(F::G)) 
\]
Composition, so defined, is associative because:
\begin{align*}
F \circ (G \circ H)
  &= shorten(T(F :: shorten(T(G ::H)) && \mbox{definition $\circ$,} \\
  &= shorten (T(F::T(G::H))) && \mbox{by lemma \ref{deferredShorteningLemma} (i),} \\
  &= shorten (T(F::G::H))    && \mbox{\parbox[t]{5cm}{see proof of associativity of composition of tight zigzags, section \ref{categorytightzigzags}.}}
\end{align*}  
and because, likewise, 
using lemma \ref{deferredShorteningLemma} (ii),
we have also that $(F \circ G) \circ H=shorten (T(F::G::H))$.

\subsection{Sequential Shortening}
\begin{lemma}
\llabel{sequentialShorteningLemma}
If $F=\zigzagTuple{f}{n}:x_1 \morph y_n$ 
is a tight stickleback in \highlight{an RR.5 range category} \catcw
and if $F$ shortens to a zigzag $\tuple{f}$ for some
morphism $f:x_1 \morph y_n$ in \catcw then  there
is a sequence of morphisms $s_1,...s_{n-1}$,
$s_i: y_i \morph y_n$ in \catcw
such that all diagrams in 
\commentary{Improve the diagramming at a later date}
\[
%\arraycolsep=12pt
\zigzagSequentialShortening{x}{y}{f}{n}{x}{g}{h}{s}
\]

commute i.e. such that
\[
f_1 \circ s_1 =f
\]
and \foreachi[n-2], 
\[
f'_i \circ s_i =  f_{i+1} \circ s_{i+1}
\] 
and
\[
f'_{n-1} \circ s_{n-1} = f_n 
\]
\end{lemma}
\begin{proof}
For any series of shortenings of $F$ to a single morphism one of this set must shorten the triangle the final zag in F i.e. 
there must be a morphism $s_{n-1}:y_{n-1} \morph y_n$
such that the triangle 
$\nudgeup{1.5cm}\downTriangleShape{x_n}{y_{n-1}}{y_n}
\downToLeft{g'}[0.6]
\downToRight{g}
\leftToRightArcing[a]{s}{30}{30} 
$ 
commutes. 
Since shortenings can be applied in any order this shortening can be applied first and the resulting zigzag 
$\tuple{f_1,f'_1,...f_{n-2}, f_{n-1} \circ s_{n-1}}$
must then shorten to the morphism $f$. 
Repeating the argument we get a morphism $s_{n-2}$
such that 
$
\begin{array}{c  c c}
&\Ob{y}{n}&               \\[0.5cm]
\Obpp{y}{n} && \Obp{y}{n} \\[0.5cm]
& \Obp{x}{n}
\end{array}
\begin{arrows}
\ncarr{ynpp}{yn}\alabel{s_{n-2}}
\ncarr{ynp}{yn}\blabel{s_{n-1}}
\ncarr{xnp}{ynpp}\alabel{f'_{n-2}}
\ncarr{xnp}{ynp}\blabel{f_{n-1}}
\end{arrows}
$
commutes and shorten $F$ to a zigzag
$\tuple{f_1,f'_1,...f_{n-3}, f_{n-2} \circ s_{n-2}}$
By iteration we arrive
at morphisms $s_1,...s_{n-1}$ which shorten $F$ to a singleton (and therefore short) zigzag $\tuple{f_1 \circ s_1}$. But since the shortening of a tight zigzag is unique 
we must have $f_1 \circ s_1=f$.
\end{proof}